\documentclass[10pt,a4paper]{article}
\usepackage{amsmath,amssymb,amsthm,graphicx}
\usepackage{hyperref}
\usepackage{enumitem}
\usepackage{mathtools}
\title{An exponential-hierarchy oracle separates worst-case derandomization from strong uniform hardness}
\author{Shuichi Hirahara, Zhenjian Lu, Hanlin Ren}
\date{doi:10.4230/LIPIcs.CCC.2023.6; CC BY 4.0}
\begin{document}
\maketitle
\noindent\textit{Editorial note.} Complete official author proof transcribed from pages22–26. Original slice lengths, machine enumeration convention, probability thresholds, heavy-query claims, bounds and final oracle-computation algorithm retained. No mathematical repair or extra assumption is supplied.\par\medskip
\noindent\textit{Allowed standard prerequisites (editorial).}\begin{itemize}
\item BPP containment in the second level of the polynomial hierarchy. Original author Theorem23 [Sipser1983] retained; its approximate probability implementation is used as the stated standard prerequisite
\end{itemize}\medskip
\paragraph{Definitions and notations (original Sections 1.1, 2.1).}
For a complexity class $\mathfrak C$, a complexity-theoretic statement is $\mathfrak C$-relativizing if it is true relative to every oracle $\mathcal O\in\mathfrak C$.
\paragraph{Definition 20.} A language $L$ is in $\mathsf{EXPH}$ if there is a constant $k$, an exponential bound $e(n)=2^{\operatorname{poly}(n)}$, and a $\operatorname{poly}(e(n))$-time algorithm $M$ such that, for every input $x\in\{0,1\}^*$,
\[x\in L\Longleftrightarrow\exists y_1\forall y_2\exists y_3\forall y_4\ldots\exists y_k,\ M(x,y_1,\ldots,y_k)=1,\]
where $y_i\in\{0,1\}^{e(n)}$ for each $i\in[k]$. Equivalently, a language is in $\mathsf{EXPH}$ if it can be decided in exponential ($2^{\operatorname{poly}(n)}$) time with access to a $\mathsf{PH}$ oracle.
\paragraph{Pseudodeterministic algorithms.} A pseudodeterministic algorithm is a probabilistic algorithm that returns a fixed output with high probability on any given input. We will also consider pseudodeterministic oracle (or advice-taking) algorithms. Such an algorithm will maintain its pseudodeterministic behavior (i.e., outputting a fixed answer with high probability for any input) when a correct oracle (or advice) is given.
\paragraph{Dense property (original Section 1.3.2).} A property $Q\subseteq\{0,1\}^*$ is dense if for every input length $n\in\mathbb N$, we have that $|Q\cap\{0,1\}^n|\geq2^n/\operatorname{poly}(n)$.
\paragraph{Theorem 23 [68].} For every oracle $\mathcal O$, $\mathsf{BPP}^{\mathcal O}\subseteq(\Sigma^p_2)^{\mathcal O}$.

\paragraph{Theorem 10.} There is an oracle $\mathcal O\in\mathsf{EXPH}$ such that
\[\mathsf{RP}^{\mathcal O}\not\subseteq\mathrm{i.o.}\,\mathsf{DTIME}^{\mathcal O}[2^n]\qquad\text{and}\qquad\mathsf{UP}^{\mathcal O}\not\subseteq\mathsf{BPTIME}^{\mathcal O}[2^n].\]

\begin{proof}[Proof of Theorem 10]
Our oracle world will consist of two oracles $P$ and $Q$ (it is easy to combine them into a single oracle). The oracle $P$ creates a hard problem in $\mathsf{RP}$ against deterministic algorithms, and the oracle $Q$ creates a hard problem in $\mathsf{UP}$ against probabilistic algorithms. In particular, every input to the oracle $P$ will be of the form $(x,r)$ where $|r|=100|x|$, and let
\[L_P:=\{x:\exists r\in\{0,1\}^{100|x|},P(x,r)=1\}.\]
We will guarantee that for every $x$, $\Pr_r[P(x,r)=1]$ is either $0$ or greater than $1/2$, thus $L_P\in\mathsf{RP}$. We also define
\[L_Q:=\{1^n:\exists r\in\{0,1\}^{100n},Q(r)=1\}.\]
We will guarantee that for every input length $100n$, $|Q\cap\{0,1\}^{100n}|\leq1$, thus $L_Q\in\mathsf{UP}$.
Let $M_1,M_2,\ldots$ be an enumeration of Turing machines running in deterministic $2^n$ time, and $N_1,N_2,\ldots$ be an enumeration of probabilistic Turing machines running in $2^n$ time (that does not necessarily satisfy the $\mathsf{BPP}$ promise).\footnote{In this paper, whenever we enumerate a list of Turing machines, we assume every machine occurs in the list infinitely many times.} We want every $M_i$ to fail to solve $L_P$ on all but finitely many input lengths, and every $N_i$ to fail to solve $L_Q$ on infinitely many input lengths. In particular, let $n_1$ be a large enough constant, and $n_i=2^{500n_{i-1}}$ for each $i\geq2$. Assuming that every machine appears in the sequence $\{N_i\}$ infinitely often, we want that for each $i\in\mathbb N$, $N_i$ fails to solve $L_Q$ on input $1^{n_i}$.
We also need the following terminologies. The $n$-th slice of $P$ is the set of inputs of the form $(x,r)$ where $|x|=n$ and $|r|=100n$. The $\ell$-th slice of $Q$ is simply the set of inputs with length $\ell$. During our construction, there will be some entries $P(x,r)$ and $Q(r)$ that are fixed, meaning that their values will never change in the subsequent construction. There will also be some inputs $x$ that are frozen, meaning that for every $r\in\{0,1\}^{100|x|}$, the values of $P(x,r)$ are fixed and will never change in the subsequent construction. If an entry $P(x,r)$ or $Q(r)$ is not fixed and, in the case of $P(x,r)$, $x$ is not frozen, then we say this entry is free.
Our construction proceeds in stages. During stage $n$, we will fix the entire $n$-th slices of $P$ and $Q$; we will possibly also fix or freeze some other entries. It will be guaranteed that before the $n$-th stage:
\begin{enumerate}
\item the number of frozen strings beyond the $n$-th slice (including the $n$-th slice) is at most $2^{n/2}$;
\item the number of fixed strings (except the frozen ones) beyond the $n$-th slice (including the $n$-th slice) is at most $2^{5n}$; and
\item each fixed or frozen entry beyond the $n$-th slice (including the $n$-th slice) is set to be $0$.
\end{enumerate}
\paragraph{Diagonalization against $M_i$.}
Let $x_1,x_2,\ldots,x_n$ be $n$ inputs of length $n$ that are not frozen. Since there are at most $2^{n/2}$ frozen inputs of length $n$, it is always possible to select $n$ unfrozen inputs. For each $i\in[n]$, we simulate the machine $M_i$ on input $x_i$. Whenever $M_i$ asks a query $P(x,r)$ or $Q(r)$, if this query is not free, then we return the corresponding value fixed before; if it is free, then we return $0$ as the answer and also fix this query to be $0$. Then $M_i$ will output an answer $\mathsf{ans}_i\in\{0,1\}$. Note that there are at most $2^{5n}+n\cdot2^n\ll2^{100n}/3$ entries in $P$ that are fixed, thus the vast majority of the entries $P(x_i,r)$ are free. We fix all these entries $P(x_i,r)$ to be $(1-\mathsf{ans}_i)$. It is easy to see that $M_i$ fails to compute $L_P$ on input $x_i$.
\paragraph{Diagonalization against $N_i$.}
If $n=10n_i$ for some $i\in\mathbb N$, then we also need to fix the $100n_i$-th slice of $Q$ so that $N_i$ fails to compute $L_Q$ on input $1^{n_i}$. (In stage $n$, this happens after the above diagonalization against $M_i$.) We simulate $N_i$ on input $1^{n_i}$. Note that $N_i$ is a probabilistic machine running in time $2^{n_i}$, thus it has $B:=2^{2^{n_i}}$ computational branches, where each branch has $1/B$ probability mass. On each branch, whenever $N_i$ asks a query $P(x,r)$ or $Q(r)$, if this query is not free, then we return the corresponding value fixed before; if it is free, then we return $0$ as the answer to the query. Note that we cannot fix the query as there are potentially as many as $B$ queries to fix. Instead, after simulating all the computational branches, we pick out the ``heavy'' queries that occur in a lot of branches and fix them:
\begin{itemize}
\item A string $x$ is heavy if at least $\frac{1}{10^4|x|^3}$ fraction of branches queried some free entry of the form $P(x,r)$.
\item An entry $Q(r)$ is heavy if at least $\frac{1}{10^4}$ fraction of branches queried $Q(r)$.
\end{itemize}
We freeze every heavy $x$ by setting every unfixed entry $P(x,r)$ to be $0$. We also fix every heavy $Q(r)$ to be $0$. Note that each computational branch of $N_i$ only makes at most $2^{n_i}$ queries. Therefore, for each input length $m$ where $n\leq m\leq2^{n_i}$, the number of heavy strings $x\in\{0,1\}^m$ is at most $10^4 2^{n_i}m^3$. The number of entries $Q(r)$ fixed before this stage is at most $2^{5n}$, and the number of heavy entries $Q(r)$ is at most $10^4 2^{n_i}$, therefore the total number of fixed entries in $Q$ after this stage is at most $2^{6n}\ll2^{100n_i}$.
After simulating every computational branch of $N_i$, let $p$ be the probability that $N_i$ outputs $1$ on input $1^n$. If $p<0.5$, then we pick a free $r\in\{0,1\}^{100n_i}$ and set $Q(r)=1$; since the number of fixed entries $Q(r)$ is $\ll2^{100n_i}$, such $r$ always exists. If $p\geq0.5$, then we do nothing.
Now we have that $N_i$ solves $L_Q$ on input $1^{n_i}$ with probability less than $2/3$ (actually, at most $1/2$). Is this always the case in the future? Consider how the future changes to the oracles might affect the behavior of $N_i$. Let $m$ be an input length where $n<m\leq2^n$. During stage $m$ (which has not happened yet), we will pick $m$ inputs $x_1,x_2,\ldots,x_m\in\{0,1\}^m$, then some entries of the form $P(x_i,r)$ will be set to be $1$. Also note that the probability mass of computational branches that each $x_i$ affects is at most $\frac{1}{10^4m^3}$. Therefore, the total probability mass of computational branches affected is at most
\[\sum_{m\in\mathbb N}\frac{m}{10^4m^3}<\frac{1}{10^3}.\]
Note that we also set at most one entry $Q(r)$ to be $1$, where $|r|=100n_i$. (The next time we diagonalize against $N_{i+1}$ happens in stage $n_{i+1}\gg2^n$.) Since this entry $Q(r)$ is not heavy, the probability mass of computational branches that it affects is at most $\frac{1}{10^4}$.
It follows that only an $\frac{11}{10^4}$ fraction of computational branches of $N_i$ will be affected in the future. That is, let $p'$ be the probability that $N_i$ outputs $1$ on input $1^{n_i}$ at the end of our construction (i.e., after stage $\gg2^n$), then $|p'-p|<\frac{11}{10^4}$. Therefore, it cannot be the case that
\[\Pr[N_i(1^{n_i})=L_Q(1^{n_i})]\geq2/3.\]
\paragraph{Clean-up.}
At the end of stage $n$, we fix every unfixed input $P(x,r)$ in the $n$-th slice to be $0$. If $n=10n_i$ then we also fix every unfixed input $Q(r)$ in the $10n$-th slice to be $0$. The number of frozen strings is at most $10^4 2^{n/10}n^3\leq2^{n/2}$ (note that at most one $N_i$ contributes to the frozen strings since the input lengths $\{n_i\}$ are very far apart). The number of fixed entries above the $(n+1)$-st slice is at most
\[2^{5n}+n\cdot2^n+\sum_{m=n}^{2^{n/10}}10^4 2^{n/10}m^3<2^{5(n+1)}.\]
\paragraph{The complexity of $\mathcal O$.}
We show that the above oracles $P$ and $Q$ are computable in $\mathsf{EXPH}$. It suffices to present a deterministic algorithm with access to a $\Sigma^p_3$ oracle that on input $1^n$, outputs the truth tables of the $n$-th slices of $P$ and $Q$ in $2^{O(n)}$ time.
We simulate each stage $n'\leq n$. Before stage $n'$, we maintain the set of fixed entries $P(x,r),Q(r)$, as well as the set of frozen strings $x$. There are at most $2^{O(n')}$ many such strings. In what follows, we use $\langle P,Q\rangle$ to denote (the length-$2^{O(n')}$ encoding of) the list of fixed entries and frozen strings.
It is easy to diagonalize against each $M_i$: we simply choose the inputs $x_1,x_2,\ldots,x_{n'}$, simulate each $M_i$ on input $x_i$, and fix the oracles accordingly. This takes $2^{O(n')}$ time.
To simulate $N_i$, we need to define the following language capturing the weight of queries. Consider simulating the probabilistic machine $N_i$ on the current oracle $\langle P,Q\rangle$. Whenever $N_i$ makes a query, if this query is not free, then we return the answer of the query; otherwise we return $0$. The following two claims are immediate corollaries of Theorem 23:
\paragraph{Claim 42.} There is a language $\mathsf{Heavy\text{-}String}(\langle P,Q\rangle,N_i,x)$ computable in $\Sigma^p_2$ such that:
\begin{itemize}
\item if at least $\frac{1.1}{10^4|x|^3}$ fraction of computational branches of $N_i$ queried some free entry of the form $P(x,-)$, then $\mathsf{Heavy\text{-}String}(\langle P,Q\rangle,N_i,x)=1$;
\item if at most $\frac{0.9}{10^4|x|^3}$ fraction of computational branches of $N_i$ queried some free entry of the form $P(x,-)$, then $\mathsf{Heavy\text{-}String}(\langle P,Q\rangle,N_i,x)=0$.
\end{itemize}
\paragraph{Claim 43.} There is a language $\mathsf{Heavy\text{-}Entry}(\langle P,Q\rangle,N_i,r)$ computable in $\Sigma^p_2$ such that:
\begin{itemize}
\item if at least $\frac{1.1}{10^4}$ fraction of branches of $N_i$ queried $Q(r)$, then $\mathsf{Heavy\text{-}Entry}(\langle P,Q\rangle,N_i,r)=1$;
\item if at most $\frac{0.9}{10^4}$ fraction of branches of $N_i$ queried $Q(r)$, then $\mathsf{Heavy\text{-}Entry}(\langle P,Q\rangle,N_i,r)=0$.
\end{itemize}
(Note that both $\mathsf{Heavy\text{-}String}$ and $\mathsf{Heavy\text{-}Entry}$ accepts inputs of length $2^{O(n')}$.)
With the aid of these two languages, it is possible to simulate stage $n'$ now. We say a string $x$ is heavy if $\mathsf{Heavy\text{-}String}(\langle P,Q\rangle,N_i,x)=1$, and an entry $Q(r)$ is heavy if $\mathsf{Heavy\text{-}Entry}(\langle P,Q\rangle,N_i,r)=1$. Since there are at most $2^{O(n')}$ heavy strings, we can find all of them in $2^{O(n')}$ time with access to an $\mathsf{NP}^{\mathsf{Heavy\text{-}String}}$ oracle. Similarly, since there are $2^{O(n')}$ heavy entries, we can find all of them in $2^{O(n')}$ time with access to an $\mathsf{NP}^{\mathsf{Heavy\text{-}Entry}}$ oracle. We freeze these heavy strings and fix these heavy entries accordingly.
Finally, we estimate the probability that $N_i$ outputs $1$ (given current versions of $P$ and $Q$) in $\Sigma^p_2$, within additive error $0.1$. If this estimation is less than $1/2$, then we pick a free $r\in\{0,1\}^{10n'}$ and set $Q(r)=1$; otherwise we do nothing.
At the end of stage $n'$, we fix every unfixed entry in the $n'$-th slice of $P(x,r)$ to be $0$. If $n=10n_i$ then we also fix every unfixed input $Q(r)$ in the $10n$-th slice to be $0$. This takes $2^{O(n')}$ time.
It is easy to see that simulating the $n$ stages takes $2^{O(n)}$ time with access to a $\Sigma^p_3$ oracle.
\end{proof}

\begin{thebibliography}{99}
\bibitem[31]{GIPS} Shafi Goldwasser, Russell Impagliazzo, Toniann Pitassi, and Rahul Santhanam. On the pseudo-deterministic query complexity of NP search problems. CCC2021, LIPIcs200,36:1--36:22. doi:10.4230/LIPIcs.CCC.2021.36.
\bibitem[33]{H18} Shuichi Hirahara. Non-black-box worst-case to average-case reductions within NP. FOCS2018,247--258. doi:10.1109/FOCS.2018.00032.
\bibitem[34]{H20} Shuichi Hirahara. Non-disjoint promise problems from meta-computational view of pseudorandom generator constructions. CCC2020,LIPIcs169,20:1--20:47. doi:10.4230/LIPIcs.CCC.2020.20.
\bibitem[68]{Sipser} Michael Sipser. A complexity theoretic approach to randomness. STOC1983,330--335. doi:10.1145/800061.808762.
\end{thebibliography}
\end{document}
